\chapter{PENUTUP}
\label{chap:penutup}

% Ubah bagian-bagian berikut dengan isi dari penutup

\section{Kesimpulan}
\label{sec:kesimpulan}

Berdasarkan eksperimen yang dilakukan pada simulasi, robot dapat berhasil menavigasi dari \emph{region} A ke \emph{region} C dengan waktu tempuh 127.26 detik. Hasil pengujian \emph{costmap} menunjukan bahwa \emph{cost scaling factor} sebesar 100 paling sesuai untuk menghasilkan jalur terpendek. Hasil ICP menunjukkan error yang wajar dan cukup akurat untuk digunakan dalam navigasi, serta dapat bekerja dengan \emph{costmap} transisi untuk menggabungkan tiap \emph{region}. Untuk lebih spesifiknya:
\begin{enumerate}
    \item Pengujian ICP pada lantai tidak rata menunjukan error lebih tinggi pada sumbu tegak lurus arah gerak dan yaw dibanding pada permukaan rate, yang disebabkan oleh perbedaan ketinggian lantai. 
    \item Obstacle lebih kompleks dan trajectory rumit di region C menghasilkan error lebih tinggi pada ICP karena robot harus menyesuaikan posisi dan orientasi lebih sering untuk menghindari rintangan.
    \item Costmap transisi dapat menggabungkan tiap \emph{region} dengan baik pada radius transisi minimal 0.26 meter.
    \item Costmap inflasi dengan \emph{cost scaling factor} 100 dapat menghasilkan jalur terpendek dengan tetap menghindari obstacle.
    \item DWA Planner dapat menghasilkan trajectory yang diikuti robot saat ada obstacle unknown
\end{enumerate}

Dengan demikian, dapat disimpulkan bahwa metode navigasi berbasis \emph{layered costmap} yang dikembangkan dalam tugas akhir ini dapat digunakan untuk menavigasi robot pada lingkungan indoor bertingkat pasca gempa bumi.

\section{Saran}
\label{chap:saran}

Beberapa saran untuk pengembangan lebih lanjut dari sistem navigasi ini
adalah sebagai berikut:

\begin{enumerate}
    \item \textbf{Pengujian di robot fisik.} Pengujian di lingkungan simulasi memberikan gambaran awal mengenai performa sistem navigasi, namun pengujian di robot fisik di lingkungan nyata akan memberikan informasi lebih mendalam mengenai tantangan yang mungkin muncul, seperti gangguan sensor, kondisi permukaan yang tidak terduga, dan interaksi dengan manusia atau objek lain.
    \item \textbf{Prekomputasi KNN individual per region.} Setiap region
    memiliki karakteristik obstacle yang berbeda, sehingga penggunaan
    tabel hash KNN terpisah per region memungkinkan penyetelan parameter
    radius dan jumlah \emph{nearest neighbors} secara independen sesuai
    kebutuhan tiap region, yang berpotensi meningkatkan akurasi dan
    efisiensi komputasi ICP.

    \item \textbf{Fusion ICP dengan IMU.} Lonjakan error ICP di zona
    transisi akibat minimnya fitur LiDAR dapat dikompensasi dengan
    sensor IMU yang menyediakan data orientasi dan percepatan pada
    frekuensi tinggi, sehingga menjaga akurasi lokalisasi secara
    konsisten di seluruh region.

    \item \textbf{Integrasi dengan Depth Camera.}
    Keterbatasan LiDAR 2D dalam mendeteksi obstacle vertikal dapat
    dilengkapi dengan data visual dari kamera atau \emph{depth sensor}
    untuk membangun representasi hibrida 3D-2D, yang memungkinkan
    robot memahami struktur lingkungan secara lebih utuh tanpa
    meninggalkan efisiensi navigasi berbasis \emph{occupancy grid}.
\end{enumerate}
